\MWChapter{Formal object and system overview}{形式对象与系统概览}
\label{ch:overview}

\section{The unified object / 统一形式对象}

The proposed semantic target is
\begin{equation}
  \mathsf{CantiluneGraph}=(\mathcal C,\mathcal R),
  \label{eq:cantilune-graph}
\end{equation}
where \(\mathcal C\) is a typed symmetric monoidal category describing legal sequential
composition \((\circ)\), parallel composition \((\otimes)\), and symmetry; \(\mathcal R\) is a
set of string-diagram rewriting rules describing identified changes to the collaboration
structure \citep{cantiluneSpec,cantiluneRFCOne}. A rule name alone is insufficient to identify an
execution: a concrete event includes the selected rule, match/redex, and derivation or
fresh-allocation data needed for replay.

\begin{figure}[H]
\centering
\resizebox{\textwidth}{!}{%
\begin{tikzpicture}[
  node distance=10mm and 10mm,
  >=Latex,
  every node/.style={font=\small\sffamily},
  layer/.style={
    rounded corners=3mm,
    draw=MWLine,
    fill=MWMist,
    minimum height=14mm,
    align=center,
    inner xsep=6mm
  },
  core/.style={
    layer,
    fill=MWLilac!9!white,
    draw=MWLilac!55,
    minimum width=48mm,
    text=MWInk
  },
  projection/.style={
    layer,
    minimum width=31mm,
    minimum height=15mm
  },
  edge/.style={->,draw=MWSlate,line width=0.8pt},
  certificate/.style={->,draw=MWCyan,line width=0.9pt,dashed}
]
  \node[layer] (declaration) {Typed generators, contracts,\\and declared policy};
  \node[core,below=of declaration] (core) {\(\mathsf{CantiluneGraph}=(\mathcal C,\mathcal R)\)\\
  legal composition + identified rewriting};

  \node[projection,below left=15mm and 31mm of core] (dag) {DAG view\\dependencies and readiness};
  \node[projection,below=15mm of core] (petri) {Petri view\\tokens and resources};
  \node[projection,below right=15mm and 31mm of core] (pi) {\(\pi\) view\\sessions and delegation};
  \node[projection,below=13mm of petri] (morphism) {Morphism view\\composition and refactoring};

  \draw[edge] (declaration) -- (core);
  \draw[edge] (core) -- (dag);
  \draw[edge] (core) -- (petri);
  \draw[edge] (core) -- (pi);
  \draw[edge] (core) -- (morphism);

  \node[font=\scriptsize\sffamily\color{MWSoftText},right=4mm of core,align=left,text width=34mm] (event) {event identity:\\[-1mm]
  rule, match, derivation, replay data};
  \draw[certificate] (event.west) -- (core.east);

  \node[
    draw=MWLilac!40,
    rounded corners=3mm,
    fit=(dag)(petri)(pi)(morphism),
    inner sep=5mm
  ] {};
\end{tikzpicture}%
}
\caption{Cantilune's intended unified semantic object. The core specifies legal typed
composition and identified rewriting; DAG, Petri, \(\pi\), and morphism are readings of that one
object. Dashed cyan denotes the certificate obligation rather than a claim that every projection
is automatically complete.}
\label{fig:architecture}
\MWSourceNote{Aligned to the formal specification and RFC-0001 \citep{cantiluneSpec,cantiluneRFCOne}.}
\end{figure}


\section{Design invariants / 设计不变量}

\begin{MWMetricGrid}
  \MWMetricCard[MWCyan]{One}{Formal object}{Structure and evolution are not separate models}
  \MWMetricCard[MWLilac]{Four}{Projection readings}{DAG, Petri, pi, and morphism views}
  \MWMetricCard[MWMint]{One}{Event identity}{Rule, match, derivation, and replay provenance}
\end{MWMetricGrid}

\begin{enumerate}
  \item \textbf{Composition is typed.} A wire connects compatible output and input ports; it is a
    contract and a typing obligation, not an untyped arrow.
  \item \textbf{Reconfiguration is explicit.} A topology, authority, or session change is proposed
    as a rule application with stated admission conditions, rather than a hidden mutation.
  \item \textbf{Views must earn agreement.} A non-identity projection requires a specified target
    transition system, event lift, preservation/reflection evidence, and terminal-observation
    conditions; monoidal structure alone is not enough.
  \item \textbf{Model work is local to a node.} The architectural proposal is: the model proposes
    within a node; the formal structure and declarative policy dispose between nodes
    \citep{cantiluneRFCOne}.
\end{enumerate}

\section{The four readings / 四个读法}

\begin{table}[H]
\centering
\caption{The four views are readings of the same intended object, not four product modules.}
\label{tab:four-readings}
\begin{tabularx}{\textwidth}{@{}L{0.20\textwidth}L{0.28\textwidth}Y@{}}
\toprule
\rowcolor{MWSoft!70}
\textbf{View} & \textbf{Question it answers} & \textbf{Formal obligation} \\
\midrule
DAG & What depends on what; what is runnable, blocked, or cyclic? & A certified observable transition reading for dependency progress \\
Petri & Which individual token, quota, lock, budget, or approval is held? & Ordered-interface, provenance-preserving resource and firing semantics \\
\(\pi\) & Which session exists; who delegated, accepted, rejected, or transferred control? & Typed/open process bridge and native late-\(\pi\) event correspondence \\
Morphism & How do pieces compose, run in parallel, and admit refactoring? & Identity reading of the symmetric monoidal structure \\
\bottomrule
\end{tabularx}
\end{table}

\begin{MWCallout}[MWAmber]{Interpretation boundary}
The diagram and table describe a formal target. They do not assert that an LLM, tool adapter,
or production control plane has already implemented these readings. The specification itself
excludes runtime, LLM, I/O, and tool/network integration from the pure formal core.
\end{MWCallout}
